% 第一章 緒論（對應英文稿 Section 1, en/front.tex）
\chapter{緒論}
\label{zh:sec:intro}

管理與資訊系統研究日益將人工智慧（artificial intelligence, AI）視為組織決策的問題。此類研究探討結合人類與 AI 決策者的決策結構 \citep{shrestha2019organizational}、管理任務之自動化與增強兩者間的相互依存 \citep{raisch2021artificial}、人類與 AI 在組織決策中的互補角色 \citep{jarrahi2018artificial}，以及 AI 的自主性、學習能力與難以解讀性所帶來的獨特管理挑戰 \citep{berente2021managing}。這些文獻皆未主張 AI 與決策之間存在結構洞（structural hole），本研究亦不將此一主張歸於它們。

此一主張其實源自某一特定類型的證據。從共現（co-occurrence）解讀結構是一種既有的做法。共詞分析（co-word analysis）依據關鍵詞的共現，描繪研究主題之間的關係或關係之闕如 \citep{callon1983translations}；文獻式知識發現（literature-based discovery）則將從未被共同引用或共同提及的文獻視為彼此無連結，並以此一闕如產生關於兩者間連結的假說 \citep{swanson1986fish}。依此類推，當 AI 功能與決策功能鮮少被同時提及時，人們很容易推論其背後的社群彼此分隔。以網路的語彙而言，此一推論指向結構洞：非冗餘聯繫對象之間的分隔，使跨越此分隔的行動者獲得資訊與控制上的利益 \citep{burt1992structural}，對管理者而言則帶來更好的構想 \citep{burt2004structural}；弱連帶則是跨越此類分隔的經典管道 \citep{granovetter1973strength}。本文所報告的研究專案將同一推論套用於 AI 與決策，而本研究即檢驗流動資料是否支持此一推論。該專案凍結的 V9 分析以文獻共現矩陣為對象，發現 16 個偏低（depleted）的 AI 功能族與決策功能族配對。專案將其解讀為 AI 與決策之間的結構性分隔，並據以建構一個多層規劃（multi-level programming, MLP）模型，描述組織如何分配資源、治理 AI，以及配置人機決策。因此，本研究所檢驗的前提，是專案本身由共現證據導出的工作前提，而非取自所引文獻的立場。

從共現推論至結構洞並不成立，而說明其原因是本研究的第一個目的。結構洞係定義於以真實資訊或資源交換關係相連之行動者的自我中心網路（ego network）上，且結構洞的測量指標各自適用於特定的資料類型（第~\ref{zh:sec:theory_holes}~節）。共現記錄的是隸屬關係，而非流動（flow）：在嚴格的二模網路（two-mode network）中，網路限制度（Burt constraint）的間接項會消失，而單模投影（one-mode projection）則會製造出封閉性（第~\ref{zh:sec:theory_cooccurrence}~節）。有效的檢驗需要帶有群組標籤之行動者之間的有向二元（dyadic）流動。專利引用長久以來被用作知識流動的證據 \citep{jaffe1993geographic}，能在母體規模上提供此類資料，且其標籤由審查員指派。本研究的第二個目的，即是在預先登錄（preregistration）之下，以此類資料檢驗該前提。

本研究探討三組研究問題。\textbf{RQ1：}AI 與決策作為知識領域，兩者之間是否存在結構洞？\textbf{RQ2：}若在領域層次並不存在，則在組織（公司）層次，於公司技術組合在機制上所隱含者之外，是否存在結構洞？\textbf{RQ3／RQ4：}根據上述答案，對於專案模型關於整合問題所在位置的設定，可得出何種推論？又應如何區辨各項競爭解釋（competing explanations）？

證據來自三條資料線（圖~\ref{zh:fig:research_design}）：文獻（研究一）、六個公開關係資料來源的篩檢，以及美國發明專利（utility patents）的母體（研究二）；在研究二中，AI 為 CPC G06N，決策（decision-making, DEC）為 CPC G06Q10/04 加上 G06Q10/06。研究二包含五項預先登錄的分析（計畫 1至計畫 5）。這些計畫是專案自有儲存庫中自行加註時間戳記的提交（commit），且計畫 2至計畫 5皆是在先前結果已知之後才撰寫的（第~\ref{zh:sec:design_prereg}~節）。研究三為一項規劃中的企業研究，不報告任何結果。

在知識流動層次，研究發現為否定性的。AI 專利與決策專利相互引用的比率約為期望值的三倍（觀察值對期望值強度〔observed-to-expected intensity, OE〕為 3.05）；兩項預先登錄的結構洞相對判準皆未成立，且判定 \texttt{HOLE\_\allowbreak NOT\_\allowbreak SUPPORTED} 在全部六種敏感度分析（sensitivity analysis）設定以及三項構念效度（construct validity）檢查（排除 G06Q10/04；匹配參照集粒度）中皆維持不變，惟部分過度連結（over-connection）可能是近乎機械性的（第~\ref{zh:sec:study2_aggregate}~節；表~\ref{zh:tab:construct}）。在公司層次，專門型（specialist）AI 公司與決策公司之間的連結少於期望，且超出專利組合構成所能解釋的分隔具有顯著性（兩種分層方式下皆為 $p = .001$）。其大小取決於所採用的尺度：在「與 1 之差距」尺度上，保留專利組合的虛無模型（portfolio-preserving null）重現了其中的 85\% 與 91\%；而在比值尺度上，專門型公司之間的連結分別為虛無模型期望值的 0.67 與 0.78，即短少 33\% 與 22\%。混合型（多角化）公司是否扮演此一分隔的中介者（broker），仍未解決。度數匹配的虛無模型顯示這類公司的限制度較低，但專利組合虛無模型並未固定度數，且沒有任何現有的虛無模型能同時固定兩者。因此，預先登錄的判定 \texttt{PARTIAL} 意指公司層次的結構洞未獲證實，但亦未被排除（第~\ref{zh:sec:study2_firm}~節）。在文獻中，16 個偏低功能配對的方向在不同隨機種子、編碼軌（tracks）與時期之間皆穩定，且在移除同詞兩端共用（same-term self-pairing）的紀錄後大致仍然成立；然而，其顯著性在實際可能的誤分類（misclassification）水準下無法維持，而模型推論的標籤仍是一項未經檢驗的非隨機誤差來源（第~\ref{zh:sec:study1}~節）。

這些結果移除了專案模型的一項正當性依據，卻未提供另一項。關於 AI 為何未能充分整合進管理決策，有數種解釋相互競爭——組織的資源分配（resource allocation）與決策權（decision rights）、資料品質、成本、法規與責任歸屬，以及任務適配性——而研究一與研究二並未偏向其中任何一項（第~\ref{zh:sec:discussion_hypothesis}~節）。MLP 模型僅將組織層面的候選解釋形式化，而研究三旨在區辨這些候選解釋，而非證實其中之一。

本研究有四項貢獻。第一，本研究說明為何從共現推論至結構洞是無效的，並在專案本身的資料中記錄四種失效模式：投影所生的團（projection cliques）、限制度在二分網路（bipartite network）上的退化（一項恆等式）、同詞兩端共用，以及模型推論的標籤。前兩者在社會網路分析中已為人所知：投影會產生封閉性 \citep{opsahl2013triadic}，投影後的中介並非二模中介 \citep{jasny2015two}，而限制度受度數與二元連帶所驅動 \citep{everett2020unpacking}。本研究並不主張發現了這些問題；本研究的貢獻在於完整追溯其對缺口推論的後果，從文獻矩陣一路延伸至以流動資料進行的檢驗。第二，本研究在五項預先登錄計畫下以有向流動資料檢驗該前提，並說明此一預先登錄能證明與不能證明之處。第三，本研究報告一項明確的否定結果。第四，本研究提出兩個陳明其限制的虛無模型（null model）：一是針對區塊結構文件的 M06-P（專利適用之保留文件結構虛無模型）——標準的交換虛無模型在此類文件上會退化——此模型並無獨立的型一誤差校準；二是保留專利組合的公司虛無模型，其並未保留公司的度數。對 MLP 研究議程的意涵係由上述結果推導而來，並非另一項研究發現。

第~\ref{zh:sec:theory}~章與第~\ref{zh:sec:design}~章闡述理論與研究設計，第~\ref{zh:sec:study1}~節與第~\ref{zh:sec:study2}~節報告研究一與研究二，第~\ref{zh:ch:discussion}~章整合各項證據，第~\ref{zh:sec:mlp}~節推導對 MLP 研究議程的意涵，第~\ref{zh:sec:limitations}~節與第~\ref{zh:sec:conclusion}~章討論研究限制並作出結論。附錄~\ref{zh:app:prereg}～\ref{zh:app:mlp} 收錄預先登錄紀錄、敏感度分析表與模型細節。
